% =====================================================================
%  SNIPPET - Lavori futuri: gate di forwarding density-aware
%
%  Da incollare nel futuro Capitolo 6 (Conclusioni e lavori futuri),
%  come sottosezione. Definisce \label{eq:fwdgate}, riutilizzabile in
%  eventuali rimandi.
%
%  Richiede amsmath nel preambolo.
% =====================================================================

\subsection{Gate di forwarding density-aware}
\label{sec:future-fwdgate}

La modalit\`a forwarded descritta nel Capitolo~\ref{ch:sistema} rilancia ogni
beacon fresco, entro TTL e non proprio, con il solo limite della coda. In reti
dense questo pu\`o produrre molti rilanci ridondanti. Un'estensione possibile \`e
un \emph{gate di forwarding} che decide, una volta sola al momento
dell'accodamento, se rilanciare un beacon ricevuto, con probabilit\`a
\begin{equation}
  P_{\text{fwd}} =
  \begin{cases}
    1 & N \le n_0,\\[2pt]
    \dfrac{n_0}{N}\cdot \max\!\big(0.3,\;1-f_q\big) & N > n_0,
  \end{cases}
  \label{eq:fwdgate}
\end{equation}
dove $N$ \`e il numero di vicini del nodo che rilancia, $n_0$ \`e un numero base
atteso di inoltranti per cascata e $f_q$ \`e il punteggio di backoff (al quadrato)
calcolato per ultimo dallo scheduler. Il fattore $n_0/N$ dirada i rilanci nelle
zone dense, cos\`i che in media solo $\approx n_0$ degli $N$ candidati rilancino.
Il fattore $\max(0.3,\,1-f_q)$ riusa lo stesso segnale di congestione locale che
governa la frequenza dei beacon propri per governare anche i rilanci: una boa che
sta gi\`a rallentando i propri beacon riduce anche i rilanci, fino a un minimo di
$0.3$.

Il gate \`e gi\`a stato progettato ma non \`e attivo nell'implementazione usata
per gli esperimenti di questa tesi. La sua valutazione --- impatto su rilanci,
collisioni, latenza e copertura, e taratura di $n_0$ --- \`e lasciata ai lavori
futuri.
